\chapter{Process\\เปลี่ยนความเชี่ยวชาญเป็นระบบที่ทำซ้ำได้}

ผู้เชี่ยวชาญมักทำหลายสิ่งได้โดยอัตโนมัติจนยากจะอธิบายว่าเหตุใดจึงได้ผล Process ทำให้ความรู้ที่ฝังอยู่ในตัวบุคคลมองเห็น ตรวจสอบ ถ่ายทอด และปรับปรุงได้ สำหรับ Tech Portfolio การอธิบาย Process มิใช่การเขียนขั้นตอนเพื่อความสวยงาม แต่เป็นหลักฐานว่าผลสำเร็จเกิดจากวิธีทำงานที่มีเหตุผล มิใช่ความบังเอิญ

\learningobjectives{
  \item แปลงวิธีทำงานที่ฝังอยู่ในประสบการณ์เป็น Process ที่ชัด
  \item ระบุ decision gate, feedback loop และเกณฑ์คุณภาพ
  \item ประเมินว่า Process ส่วนใดควรถ่ายทอดหรือปรับปรุง
}

\section{Process ที่มีความหมาย}
รายการกิจกรรมบอกว่าเกิดอะไรขึ้น แต่ Process อธิบายความสัมพันธ์ระหว่างกิจกรรม ข้อมูล และการตัดสินใจ Process ที่ดีจึงประกอบด้วย input, ขั้นตอนหลัก, จุดตัดสินใจ, output, ผู้รับผิดชอบ และเกณฑ์ที่บอกว่าจะเดินหน้าหรือย้อนกลับ

\begin{longtable}{P{.18\textwidth}P{.24\textwidth}P{.45\textwidth}}
\toprule
\textbf{ระยะ} & \textbf{เป้าหมาย} & \textbf{ตัวอย่างเกณฑ์ตัดสินใจ}\\
\midrule
Discover & เข้าใจปัญหาและผู้ใช้ & มีหลักฐานจากผู้ใช้หลายกลุ่มว่ายืนยัน pain point เดียวกัน\\
Define & เลือกปัญหาที่คุ้มค่าแก่การแก้ & ผลกระทบสูงและอยู่ในขอบเขตความสามารถของทีม\\
Prototype & ทดสอบสมมติฐานด้วยต้นทุนต่ำ & ต้นแบบตอบคำถามสำคัญโดยไม่ต้องสร้างระบบเต็ม\\
Pilot & ตรวจสอบในบริบทจริง & ผู้ใช้ทำงานได้จริงและผลลัพธ์ผ่านเกณฑ์ขั้นต่ำ\\
Scale/Learn & ขยายและปรับระบบ & คุณภาพคงที่ มีเจ้าของกระบวนการและ feedback loop\\
\bottomrule
\end{longtable}

\section{Decision Gate และ Feedback Loop}
จุดแข็งของ Process ไม่ได้อยู่ที่เดินไปข้างหน้าเสมอ แต่อยู่ที่รู้ว่าเมื่อใดควรหยุด เปลี่ยน หรือย้อนกลับ Decision gate เช่น \enquote{จะเข้าสู่ pilot ต่อเมื่อผู้ใช้เป้าหมายอย่างน้อย 5 คนทำภารกิจสำเร็จโดยไม่ต้องมีทีมช่วย} ทำให้การตัดสินใจตรวจสอบได้ ส่วน feedback loop ทำให้ข้อมูลจากผลลัพธ์ย้อนมาปรับวิธีทำงาน

\begin{examplebox}
ทีมหนึ่งเคยใช้เวลาสามเดือนสร้างต้นแบบก่อนพบว่าผู้ใช้ไม่ต้องการ หลังถอดบทเรียน ทีมกำหนด Process ใหม่: สัมภาษณ์ 12 ราย สรุปสมมติฐาน 3 ข้อ สร้างต้นแบบกระดาษ และทดสอบภายในสองสัปดาห์ หากผู้ใช้ไม่สามารถอธิบายคุณค่าได้อย่างน้อย 8 ราย ทีมจะกลับไปนิยามปัญหาใหม่ Process นี้ลดต้นทุนการเรียนรู้และสร้างหลักฐานเร็วขึ้น
\end{examplebox}

\section{ถ่ายทอดโดยไม่ทำให้แข็งตัว}
การทำ Process ให้ชัดไม่ได้หมายความว่าทุกกรณีต้องทำเหมือนกัน Mentor ควรช่วย Mentee แยก \textbf{หลักการที่ต้องรักษา} ออกจาก \textbf{วิธีที่ปรับตามบริบทได้} เช่น หลักการคือทดสอบกับผู้ใช้ก่อนลงทุนสูง ส่วนจำนวนผู้ใช้หรือรูปแบบต้นแบบอาจเปลี่ยนได้

\diagnostic{ส่วนใดของงานที่ยังทำได้ดีเฉพาะเมื่อ Mentee อยู่ในห้อง และความรู้ใดต้องทำให้มองเห็นเพื่อให้ทีมทำต่อได้}

\begin{mentornote}
เริ่มถอด Process จากเหตุการณ์จริงหนึ่งครั้ง ไม่ใช่จาก Process ในอุดมคติ ความแตกต่างระหว่าง \enquote{สิ่งที่ควรเกิด} กับ \enquote{สิ่งที่เกิดจริง} มักเผยจุดปรับปรุงที่สำคัญ
\end{mentornote}

\begin{keytakeaway}
Process แปลงความเชี่ยวชาญส่วนบุคคลเป็นความสามารถของทีม โดยทำให้เหตุผล จุดตัดสินใจ เกณฑ์คุณภาพ และวงจรเรียนรู้มองเห็นได้
\end{keytakeaway}

\begin{reflection}
วาด Process ของงานสำคัญหนึ่งชิ้นตั้งแต่รับโจทย์จนส่งมอบ ใส่จุดตัดสินใจอย่างน้อยสามจุด ระบุข้อมูลที่ใช้ และทำเครื่องหมายส่วนที่ยังพึ่งความรู้ในตัวบุคคลมากเกินไป
\blankline\blankline
\end{reflection}

\chaptersummary{Process ที่ดีมิใช่เพียง workflow แต่เป็นตรรกะของการทำงานที่ตรวจสอบและเรียนรู้ได้ การทำให้ Process ชัดช่วยลดความเสี่ยง ถ่ายทอดความสามารถ และทำให้ผลลัพธ์น่าเชื่อถือขึ้น}
